\documentclass[10pt,a4paper,ssfamily]{exam}
\usepackage[brazil]{babel}
\usepackage[numbers,sort&compress]{natbib}
\usepackage[utf8]{inputenc}
\renewcommand{\baselinestretch}{1.0}
\usepackage{epsfig}
\usepackage{mhchem}
\usepackage{graphicx}
\usepackage{makeidx}
\usepackage{multicol}
\usepackage{multirow}
\usepackage{amssymb}
\usepackage{amsmath}
\DeclareMathOperator{\sen}{sen}
\newcommand{\listfont}{\bf\fontencoding{OT1}\sffamily\large}
\newcommand{\titlesfont}{\fontencoding{OT1}\sffamily\large}
\renewcommand{\baselinestretch}{1.0}
\newcommand{\Mg}{Mg$^{2+}$}
\renewcommand{\t}{\textrm}
\newcommand{\cc}[1]{[\textrm{#1}]}
\usepackage{enumitem}
\newcommand{\bmat}{\left[\begin{array}{cc}}
\newcommand{\emat}{\end{array}\right]}
\newcommand{\nomera}{
  \noindent
  \begin{tabular}{|p{0.7\textwidth}|p{0.3\textwidth}|}
  Nome: & RA: \\
  \hline
  \end{tabular}
}

\newcommand{\qini}{\begin{enumerate}[resume]\item}
\newcommand{\qend}{\end{enumerate}}

\newcommand{\oC}{$^\circ$C}
\newcommand{\kcalmol}{~kcal~mol$^{-1}$}
\newcommand{\moll}{~mol~L$^{-1}$}
\newcommand{\molkg}{~mol~kg$^{-1}$}
\newcommand{\gl}{~g~L$^{-1}$}
\newcommand{\e}[1]{$\times 10^{#1}$}

\newcommand{\eps}{\varepsilon}

\newcommand{\Resposta}[1]{\vspace{0.5cm}\noindent{\bf Resposta:}\\ #1
\begin{center}\rule{\textwidth}{0.4pt}\end{center}}
\renewcommand{\Resposta}[1]{}
\newcommand{\resposta}[1]{\vspace{0.5cm}\noindent{\bf Resposta:}\\ #1
\begin{center}\rule{\textwidth}{0.4pt}\end{center}}


\begin{document}
\sffamily
\pagestyle{empty}
\nomera

\begin{center}

{\bf QG101 - Química Geral}

{\bf Aula 3 - Estrutura Atômica}\\
\underline{Justifique suas respostas.}
\end{center}

\begin{center}{\bf Questões}\end{center}

% ---------------------------------------------------------------
\qini
O \textbf{cálcio} (\ce{Ca}) é essencial para a estrutura óssea e para
diversas aplicações em engenharia civil (cal, cimento). Seu isótopo mais
abundante é o ${}^{40}_{20}\text{Ca}$.

\begin{enumerate}
  \item Determine o número de prótons, nêutrons e elétrons do átomo
        neutro de ${}^{40}_{20}\text{Ca}$.
  \item Escreva a configuração eletrônica completa do \ce{Ca} e sua
        forma abreviada usando o gás nobre anterior como núcleo.
  \item A qual grupo e período da tabela periódica o Ca pertence?
        Quantos elétrons de valência possui? Com base nisso, qual íon
        o Ca tende a formar?
\end{enumerate}

\vspace{2cm}
\qend

% ---------------------------------------------------------------
\qini
O \textbf{cloro} (\ce{Cl}) é amplamente usado no tratamento de água
potável. Ocorre naturalmente como mistura de dois isótopos estáveis:

\begin{center}
\begin{tabular}{|c|c|c|}
\hline
Isótopo & Massa (u) & Abundância natural \\
\hline
${}^{35}\text{Cl}$ & $34{,}969$ & $75{,}77\%$ \\
${}^{37}\text{Cl}$ & $36{,}966$ & $24{,}23\%$ \\
\hline
\end{tabular}
\end{center}

\begin{enumerate}
  \item Calcule a massa atômica média do cloro.
        Compare com o valor tabelado ($M = 35{,}45~\text{u}$).
  \item O ${}^{35}\text{Cl}$ e o ${}^{37}\text{Cl}$ são isótopos?
        Justifique com base na definição de número atômico e número
        de massa. Quantos nêutrons tem cada um?
  \item Escreva a configuração eletrônica do \ce{Cl} ($Z = 17$) e
        indique quantos elétrons de valência ele possui.
        O \ce{Cl} tende a ganhar ou a perder elétrons? Qual íon forma?
\end{enumerate}

\vspace{2cm}
\qend

% ---------------------------------------------------------------
\qini
O \textbf{ferro} (\ce{Fe}, $Z = 26$) é o metal estrutural mais
utilizado na engenharia. Sua tendência à corrosão está diretamente
ligada à configuração eletrônica e à facilidade de ionização.

\begin{enumerate}
  \item Escreva a configuração eletrônica completa do \ce{Fe} e
        sua forma abreviada com gás nobre.
  \item O \ce{Fe} pode se oxidar formando \ce{Fe^{2+}} (perde 2
        elétrons) ou \ce{Fe^{3+}} (perde 3 elétrons). Escreva a
        configuração eletrônica de cada íon. De qual subnível
        são retirados os elétrons em cada caso? Justifique.
  \item Entre o \ce{Fe} ($Z=26$) e o \ce{Cu} ($Z=29$), qual possui
        maior raio atômico? Justifique com base na tendência de
        raio atômico ao longo de um período.
\end{enumerate}

\vspace{2cm}
\qend

\end{document}
